\begin{trapbox}
\begin{description}[style=nextline, leftmargin=2.8em, labelsep=0.5em, itemsep=0.35em]
    \item[\textbf{\textcolor{CTrap}{易错点一（忽略正实性）}}] 认为不等式对所有实数成立，忽略 $a,b,c>0$ 的条件。
    \item[\textbf{\textcolor{CTrap}{正确理解（严格正数）：}}] 必须确保 $a,b,c$ 均为正实数，否则不等式可能失效。
    \item[\textbf{\textcolor{CTrap}{错因分析（符号影响）：}}] 若变量为负，分母可能为负导致不等号方向改变，或分母为零无定义。
\end{description}

\medskip
\begin{description}[style=nextline, leftmargin=2.8em, labelsep=0.5em, itemsep=0.35em]
    \item[\textbf{\textcolor{CTrap}{易错点二（等号条件错）}}] 误以为等号在 $a,b,c$ 满足某种比例时成立，而非全相等。
    \item[\textbf{\textcolor{CTrap}{正确理解（全等取等）：}}] 等号成立当且仅当 $a=b=c$，没有其他情况。
    \item[\textbf{\textcolor{CTrap}{错因分析（对称性误解）：}}] 由于不等式对称，但最小值点唯一，误以为对称性允许更多取等点。
\end{description}

\medskip
\begin{description}[style=nextline, leftmargin=2.8em, labelsep=0.5em, itemsep=0.35em]
    \item[\textbf{\textcolor{CTrap}{易错点三（推广形式误用）}}] 在使用推广形式时，错误地设 $s$ 或忘记分母是 $s-a$。
    \item[\textbf{\textcolor{CTrap}{正确理解（和表示分母）：}}] 推广形式中 $s=a+b+c$，分母应为 $s-a = b+c$ 等。
    \item[\textbf{\textcolor{CTrap}{错因分析（记号混淆）：}}] 混淆 $s$ 的定义，导致代数变形错误。
\end{description}
\end{trapbox}
